% =============================================================
%  CS102-F26 Digital Design and Computer Architecture I
%  Lecture 22 -- Building and Programming the DX8 (Unit 22)
%  Follows Section 6.6.4 and the Epilogue of
%  A Student's Computer Architecture Book I, and the DX8 circuit of its
%  Figure 6.13 (chapter_6.circ of the DX8 course material).
%
%  Build the slides and the print handout:  npm run materials -- lecture-22
% =============================================================
\makeatletter
\def\input@path{{templates/latex/}{../}}
\makeatother

\documentclass[10pt]{wosm-slides}
\usepackage{cs102-lectures}[2026/10/20]

\title[Lecture 22 \textperiodcentered\ Building and Programming the DX8]{Building and
  Programming the DX8}
\subtitle{Lecture 22: putting it all together}

% ---------- Helpers for this lecture --------------------------
\tikzset{
  part/.style={draw=wosm-text, line width=0.6pt, fill=wosm-bg, align=center,
    font=\scriptsize, minimum width=1.5cm, minimum height=0.75cm, inner sep=2pt},
  group/.style={draw=wosm-muted, line width=0.5pt, dashed},
  grouplabel/.style={font=\scriptsize, text=wosm-muted},
  hot/.style={draw=wosm-brick, line width=1.5pt},
}

% The DX8 as drawn in Lecture 21 (the book's Figures 6.8 to 6.11).
\newcommand*\datapath[1]{%
  \begin{tikzpicture}[x=0.92cm, y=0.92cm,
      ip-ram/.style={}, ram-rf/.style={}, rf-ir/.style={}, rf-ab/.style={},
      ra-alu/.style={}, rb-alu/.style={}, ir-alu/.style={}, st-ctl/.style={},
      ctl-alu/.style={}, alu-out/.style={}, alu-flag/.style={}, out-rf/.style={},
      flag-rf/.style={}, ir-ctl/.style={}, ctl-ip/.style={}, #1]
    % groups
    \draw[group] (1.65,-1.85) rectangle (3.55,1.75);
    \node[grouplabel, anchor=north] at (2.6,-1.85) {memory unit};
    \draw[group] (4.0,-2.85) rectangle (11.45,2.4);
    \node[grouplabel, anchor=south east] at (11.45,-2.85) {CPU};
    % parts
    \node[part, minimum width=1.4cm] (in) at (0,0) {input unit\\ + buttons};
    \node[part, minimum width=1.4cm, minimum height=0.8cm] (ram) at (2.6,1.0) {RAM};
    \node[part, minimum width=1.4cm] (ip) at (2.6,-1.05) {instruction\\ pointer (IP)};
    \node[part, minimum height=0.8cm] (rf) at (5.3,1.0) {register file\\ R0--R7};
    \node[part] (ir) at (5.3,-0.3) {instruction\\ register (IR)};
    \node[part, minimum width=0.9cm, minimum height=0.5cm] (ra) at (7.2,1.55) {regA};
    \node[part, minimum width=0.9cm, minimum height=0.5cm] (rb) at (8.4,1.55) {regB};
    \node[part, minimum width=2.2cm, minimum height=0.9cm] (alu) at (7.8,0.1) {ALU};
    \node[part, minimum width=1.2cm, minimum height=0.45cm] (out) at (7.6,-1.0) {regOUT};
    \node[part, minimum width=1.2cm, minimum height=0.45cm] (flag) at (7.6,-1.7) {regFLAG};
    \node[part, minimum width=1.4cm] (st) at (10.4,1.55) {stepper};
    \node[part, minimum width=1.4cm] (ctl) at (10.4,0.1) {controller};
    % input
    \draw[wire] (in.east) -- (1.1,0);
    \draw[arrow] (1.1,0) |- (ram.west);
    \draw[arrow] (1.1,0) |- (ip.west);
    \fill[wosm-text] (1.1,0) circle (1.5pt);
    % paths
    \draw[arrow, ip-ram] (ip) -- (ram);
    \draw[arrow, {Stealth[length=5pt]}-{Stealth[length=5pt]}, ram-rf]
      (ram.east |- 0,1.15) -- (rf.west |- 0,1.15);
    \draw[arrow, rf-ir] (rf) -- (ir);
    \draw[wire, rf-ab] (rf.east) -- ++(0.25,0) |- (8.4,2.1);
    \draw[arrow, rf-ab] (7.2,2.1) -- (ra);
    \draw[arrow, rf-ab] (8.4,2.1) -- (rb);
    \fill[wosm-text] (7.2,2.1) circle (1.5pt);
    \draw[arrow, ra-alu] (ra) -- (ra |- alu.north);
    \draw[arrow, rb-alu] (rb) -- (rb |- alu.north);
    \draw[arrow, ir-alu] (ir.east) -- (ir.east -| alu.west);
    \draw[arrow, st-ctl] (st) -- (ctl);
    \draw[arrow, ctl-alu] (ctl) -- (alu);
    \draw[arrow, alu-out] (alu.south -| out) -- (out);
    \draw[arrow, alu-flag] (8.7,-0.39) |- (flag.east);
    \draw[wire, flag-rf] (flag.west) -| (4.25,-1.0);
    \draw[arrow, out-rf] (out.west) -| (4.25,0.85) -- (rf.west |- 0,0.85);
    \fill[wosm-text] (4.25,-1.0) circle (1.5pt);
    \draw[arrow, ir-ctl] (ir.south) |- (10.1,-2.35) -- (10.1,0 |- ctl.south);
    \draw[arrow, ctl-ip] (10.75,0 |- ctl.south) |- (2.6,-2.6) -- (ip.south);
  \end{tikzpicture}}

% =============================================================
\begin{document}

\maketitle

\begin{frame}{Today}
  \tableofcontents
\end{frame}

% =============================================================
\section{The whole DX8}

\begin{frame}{All the parts, together}
  \begin{center}
    \datapath{}
  \end{center}
  \begin{itemize}
    \item The book's Figure 6.13 connects everything of Units 15 to 21, with a few
      more registers to hold the instruction, the ALU's inputs and its results.
  \end{itemize}
\end{frame}

\begin{frame}{One instruction in eight steps}
  \begin{center}
    \small
    \begin{tabular}{@{}cll@{}}
      \toprule
      Step & Stage & In the book's circuit \\
      \midrule
      $s_1$ & fetch   & the byte at address IP, $b_{15:8}$, is copied into R0 \\
      $s_2$ & fetch   & the byte at IP $+ 1$, $b_{7:0}$, is copied into R1 \\
      $s_3$ & decode  & R0 goes to the instruction register: $b_{15}$, opcode, register \\
      $s_4$ & decode  & R1 goes to the argument register: \texttt{val} or the address \\
      $s_5$ & execute & the ALU's result goes to regOUT, its flags to regFLAG \\
      $s_6$ & execute & (nothing more to do) \\
      $s_7$ & load    & results are written back: register file, regA, regB or RAM \\
      $s_8$ & load    & the IP moves on to IP $+ 2$, or to \texttt{val} after a jump \\
      \bottomrule
    \end{tabular}
  \end{center}

  \begin{itemize}
    \item regOUT and regFLAG are loaded at step $s_5$ of \emph{every} instruction, so
      they are only temporary. A result that must last is written to a register in
      step $s_7$.
  \end{itemize}
\end{frame}

\begin{frame}{The load stage, instruction by instruction}
  \begin{center}
    \begin{tabular}{@{}lll@{}}
      \toprule
      Instruction & In step $s_7$ & In step $s_8$ \\
      \midrule
      ALU instruction & regOUT $\to$ the register in $b_{10:8}$ & IP $+ 2$ \\
      \texttt{move} & \texttt{val} $\to$ the register in $b_{10:8}$ & IP $+ 2$ \\
      \texttt{arg1}, \texttt{arg2} & the register in $b_{10:8}$ $\to$ regA, regB
        & IP $+ 2$ \\
      \texttt{store} & the register in $b_{10:8}$ $\to$ $M[\texttt{memAddr}]$
        & IP $+ 2$ \\
      \texttt{jmp} & & \texttt{val} \\
      \texttt{jmpgt}, \ldots & & \texttt{val} if the flag is set, else IP $+ 2$ \\
      \bottomrule
    \end{tabular}
  \end{center}

  \begin{itemize}
    \item An ALU instruction does use its register field: it names the register that
      receives the result. \texttt{cmp} writes 0 there, since its answer is in the
      flags.
    \item Each instruction is 2 bytes long, so ``the next instruction'' (regIP $+ 1$
      in Table 6.2) is 2 addresses further on.
  \end{itemize}
\end{frame}

\begin{frame}{Rules for DX8 programs}
  \begin{itemize}
    \item \textbf{Two bytes per instruction}: $b_{15:8}$ at an even address, $b_{7:0}$
      at the next one.
    \item \textbf{Start at 0}: reset sets the IP to 0 and the stepper to step 1.
    \item \textbf{Leave R0 and R1 alone}: every fetch writes the instruction into
      them. Keep your values in R2 to R7.
    \item \textbf{Choose where results go}: an ALU instruction writes its result into
      the register in $b_{10:8}$.
    \item \textbf{The flags compare regA with regB}, as they are at step $s_5$ of the
      instruction: a conditional jump tests the values put there by the last
      \texttt{arg1} and \texttt{arg2}.
    \item \textbf{There is no halt}: end with a jump to the jump itself.
  \end{itemize}
\end{frame}

% =============================================================
\section{Programming the DX8}

\begin{frame}{Switches and buttons}
  \begin{columns}[T, totalwidth=\textwidth]
    \column{0.48\textwidth}
      \begin{definition}[Button]
        A (push) \emph{button} sends 1 as long as it is pressed; its output drops to
        0 as soon as you lift your finger.
      \end{definition}
    \column{0.48\textwidth}
      \begin{definition}[Switch]
        A \emph{switch} sends 1 once you turn it on, and keeps sending 1 until you
        turn it off again.
      \end{definition}
  \end{columns}

  \medskip
  The DX8's ``keyboard'' \unitref{15}:
  \begin{itemize}
    \item 8 \emph{data} switches and 8 \emph{address} switches;
    \item a \emph{write} button, which stores the data at the address;
    \item a \emph{reset} switch.
  \end{itemize}
\end{frame}

\begin{frame}{Typing in a program}
  \begin{enumerate}
    \item Turn \alert{on} the reset switch.
    \item Set the data switches to the next byte of the program.
    \item Set the address switches to its address.
    \item Push the write button (as if you turned a switch on and off at once).
    \item Repeat 2 to 4 until every byte of the program and its data is in the RAM.
    \item Turn \alert{off} the reset switch: the DX8 starts running from address 0.
  \end{enumerate}

  \begin{remark}
    The program is just bytes in the RAM, typed in like data: the Von Neumann idea
    \unitref{15}, at work.
  \end{remark}
\end{frame}

\begin{frame}{A program: $3 \times 4$ by repeated addition}
  \begin{columns}[T, totalwidth=\textwidth]
    \column{0.49\textwidth}
      \footnotesize
      \begin{tabular}{@{}lll@{}}
        \toprule
        Addr. & Bytes & Instruction \\
        \midrule
        \texttt{00} & \texttt{12 00} & \texttt{move} R2, 0 \quad (the product) \\
        \texttt{02} & \texttt{13 04} & \texttt{move} R3, 4 \\
        \texttt{04} & \texttt{14 03} & \texttt{move} R4, 3 \quad (the counter) \\
        \texttt{06} & \texttt{15 FF} & \texttt{move} R5, $-1$ \\
        \texttt{08} & \texttt{16 00} & \texttt{move} R6, 0 \\
        \texttt{0A} & \texttt{1A 00} & \texttt{arg1} R2 \quad $\leftarrow$ loop \\
        \texttt{0C} & \texttt{23 00} & \texttt{arg2} R3 \\
        \texttt{0E} & \texttt{8A 00} & \texttt{add} $\to$ R2 \quad (R2 $\leftarrow$ R2 $+$ 4) \\
        \texttt{10} & \texttt{1C 00} & \texttt{arg1} R4 \\
        \texttt{12} & \texttt{25 00} & \texttt{arg2} R5 \\
        \bottomrule
      \end{tabular}
    \column{0.49\textwidth}
      \footnotesize
      \begin{tabular}{@{}lll@{}}
        \toprule
        Addr. & Bytes & Instruction \\
        \midrule
        \texttt{14} & \texttt{8C 00} & \texttt{add} $\to$ R4 \quad (R4 $\leftarrow$ R4 $-$ 1) \\
        \texttt{16} & \texttt{1C 00} & \texttt{arg1} R4 \\
        \texttt{18} & \texttt{26 00} & \texttt{arg2} R6 \\
        \texttt{1A} & \texttt{87 00} & \texttt{cmp} $\to$ R7 \quad (R4 vs 0) \\
        \texttt{1C} & \texttt{30 0A} & \texttt{jmpgt} \texttt{0A} \\
        \texttt{1E} & \texttt{1A 00} & \texttt{arg1} R2 \\
        \texttt{20} & \texttt{26 00} & \texttt{arg2} R6 \\
        \texttt{22} & \texttt{8F 00} & \texttt{add} $\to$ R7 \quad (R2 $+$ 0) \\
        \texttt{24} & \texttt{28 24} & \texttt{jmp} \texttt{24} \quad (stop) \\
        \bottomrule
      \end{tabular}
  \end{columns}

  \medskip
  \small Addresses and bytes in hexadecimal. $-1$ is \texttt{FF} in two's complement
  \unitref{16}, so adding it counts down.
\end{frame}

\begin{frame}{Running it}
  \begin{columns}[c, totalwidth=\textwidth]
    \column{0.42\textwidth}
      \centering
      \begin{tabular}{@{}lccc@{}}
        \toprule
        After & R2 & R4 & jump? \\
        \midrule
        set-up & 0 & 3 & \\
        pass 1 & 4 & 2 & yes \\
        pass 2 & 8 & 1 & yes \\
        pass 3 & 12 & 0 & no \\
        \bottomrule
      \end{tabular}
    \column{0.54\textwidth}
      \begin{itemize}
        \item Each pass adds 4 to R2 and takes 1 from R4. While $R4 > 0$,
          \texttt{jmpgt} goes back to address \texttt{0A}.
        \item At the end, regOUT shows 12 (\texttt{0000\,1100}), and the DX8 keeps
          jumping to \texttt{24}.
        \item 38 instructions of 8 cycles each run before the stop: 304 clock cycles.
      \end{itemize}
  \end{columns}

  \begin{remark}
    The DX8 has no multiplication instruction; loops of additions do the job. Real
    processors often have a multiplier circuit for it.
  \end{remark}
\end{frame}

% =============================================================
\section{The project}

\begin{frame}{The DX8 project}
  \begin{itemize}
    \item The circuit-design project of this module: build the DX8 computer, or an
      extension of it, in Logisim-evolution, and run your own programs on it. It
      counts for 40\% of the module.
    \item Build it from your own parts of Units 6 to 21, and test each one before
      you connect it.
    \item Ideas for an extension: a \texttt{sub} instruction in one of the 12 unused
      opcodes; a \texttt{halt}; an output display for regOUT; a \texttt{load} of a
      byte from the RAM into a register.
  \end{itemize}

  \begin{remark}
    The book's epilogue admits that the DX8 was built without the setup and hold
    checks of Unit 14. Doing them for your DX8 is a good extension too.
  \end{remark}
\end{frame}

\begin{frame}{Try it}
  \begin{enumerate}
    \item In which step does the IP change, and by how much?\par
      \uncover<2->{\small In $s_8$: by 2, or to \texttt{val} after a jump.}
    \item Write \texttt{move} R5, \texttt{0x0A} as two bytes.\par
      \uncover<3->{\small \texttt{15 0A}: $0\,0010\,101$, then $0000\,1010$.}
    \item Which one byte of the program would you change to compute $3 \times 5$?\par
      \uncover<4->{\small The byte at address \texttt{03}: \texttt{04} becomes
        \texttt{05}.}
    \item Why may a program not keep a value in R1?\par
      \uncover<5->{\small Every fetch writes the second byte of the instruction into
        R1.}
    \item How many clock cycles does a program of 10 instructions take on the DX8,
      without jumps?\par
      \uncover<6->{\small $10 \cdot 8 = 80$.}
  \end{enumerate}
\end{frame}

% =============================================================
\section{Summary}

\begin{frame}{Summary}
  \begin{itemize}
    \item The DX8 runs every instruction in 8 steps: two bytes are fetched into R0 and
      R1, decoded into the instruction and argument registers, executed by the ALU,
      and written back; then the IP moves on.
    \item Results last only if written back: an ALU instruction writes into the
      register in $b_{10:8}$.
    \item Programs start at address 0, use R2 to R7, and end with a jump to itself.
    \item To program the DX8: reset on, type each byte with the data and address
      switches and the write button, reset off.
  \end{itemize}
\end{frame}

\begin{frame}{Before the last lecture}
  \begin{itemize}
    \item Read Section~6.6.4 and the Epilogue of the book.
    \item Put your DX8 together in Logisim-evolution, type in the $3 \times 4$
      program, and check that regOUT shows 12.
    \item Write a program of your own: for example, the larger of two numbers, or
      $1 + 2 + \cdots + 10$.
  \end{itemize}

  \begin{block}{Next lecture \unitref{23}}
    Course review, and the exam.
  \end{block}
\end{frame}

\end{document}
